%% Le lingue utilizzate, che verranno passate come opzioni al pacchetto babel. Come sempre, l'ultima indicata sarà quella primaria.
%% Se si utilizzano una o più lingue diverse da "italian" o "english", leggere le istruzioni in fondo.
\def\thudbabelopt{english}
%% Valori ammessi per target: bach (tesi triennale), mst (tesi magistrale), phd (tesi di dottorato).
%% Valori ammessi per aauheader: '' (vuoto -> nessun header Alpen Adria Univeristat), aics (Department of Artificial Intelligence and Cybersecurity), informatics (Department of Informatics Systems). Il nome del dipartimento è allineato con la versione inglese del logo UniUD.
%% Valori ammessi per style: '' (vuoto -> stile moderno), old (stile tradizionale).
\documentclass[target=mst,aauheader=aics,style=]{thud}

%% --- Informazioni sulla tesi ---
%% Per tutti i tipi di tesi
% Scommentare quello di interesse, o mettete quello che vi pare
\course{Artificial Intelligence \& Cybersecurity}
%\course{Internet of Things, Big Data e Web}
%\course{Matematica}
%\course{Comunicazione Multimediale e Tecnologie dell'Informazione}
\title{Titolo da vedere}
\author{Nicola Flego}
\supervisor{Prof.\ Giuseppe Serra}
\cosupervisor{Arch.\ Rambaldo Melandri \and Dott.\ Giorgio Perozzi}
\tutor{Guido Necchi}
%% Campi obbligatori: \title, \author e \course.
%% Altri campi disponibili: \reviewer, \tutor, \chair, \date (anno accademico, calcolato in automatico), \rights
%% Con \supervisor, \cosupervisor, \reviewer e \tutor si possono indicare più nomi separati da \and.


%% --- Pacchetti consigliati ---
%% pdfx: per generare il PDF/A per l'archiviazione. Necessario solo per la versione finale
\usepackage[a-1b]{pdfx}
%% hyperref: Regola le impostazioni della creazione del PDF... più tante altre cose. Ricordarsi di usare l'opzione pdfa.
\usepackage[pdfa]{hyperref}
%% tocbibind: Inserisce nell'indice anche la lista delle figure, la bibliografia, ecc.
\usepackage{tocbibind}
\usepackage{amsmath}
\usepackage{enumitem}
\usepackage{amssymb}

%% --- Stili di pagina disponibili (comando \pagestyle) ---
%% sfbig (predefinito): Apertura delle parti e dei capitoli col numero grande; titoli delle parti e dei capitoli e intestazioni di pagina in sans serif.
%% big: Come "sfbig", solo serif.
%% plain: Apertura delle parti e dei capitoli tradizionali di LaTeX; intestazioni di pagina come "big".

\begin{document}
\maketitle

%% Dedica (opzionale)
\begin{dedication}
	Al mio cane,\par per avermi ascoltato mentre ripassavo le lezioni.
\end{dedication}

%% Ringraziamenti (opzionali)
\acknowledgements
Sed vel lorem a arcu faucibus aliquet eu semper tortor. Aliquam dolor lacus, semper vitae ligula sed, blandit iaculis leo. Nam pharetra lobortis leo nec auctor. Pellentesque habitant morbi tristique senectus et netus et malesuada fames ac turpis egestas. Fusce ac risus pulvinar, congue eros non, interdum metus. Mauris tincidunt neque et aliquam imperdiet. Aenean ac tellus id nibh pellentesque pulvinar ut eu lacus. Proin tempor facilisis tortor, et hendrerit purus commodo laoreet. Quisque sed augue id ligula consectetur adipiscing. Vestibulum libero metus, lacinia ac vestibulum eu, varius non arcu. Nam et gravida velit.

%% Sommario (opzionale)
\abstract
Nunc ac dignissim ipsum, quis pulvinar elit. Mauris congue nec leo ornare lobortis. Nulla hendrerit pretium diam nec lobortis. Nullam aliquam laoreet nisl, sit amet facilisis lectus accumsan ut. Duis et elit hendrerit metus venenatis condimentum. Integer id eros molestie, interdum leo sit amet, aliquet metus. Integer fermentum tristique magna, vel luctus neque rhoncus vel. Ut hendrerit et quam et semper. Mauris egestas, odio sed aliquet luctus, magna orci euismod odio, vitae lacinia tellus tellus non lectus. Aliquam urna neque, porta et mattis aliquam, congue sit amet lorem. In ultrices augue sit amet ante vehicula, vitae rhoncus turpis auctor. Donec porta scelerisque eros, at mollis enim imperdiet ut. 

%% Indice
\tableofcontents

%% Lista delle tabelle (se presenti)
\listoftables

%% Lista delle figure (se presenti)
\listoffigures

%% Corpo principale del documento
\mainmatter

%% Parte
%% La suddivisione in parti è opzionale; solitamente sono sufficienti i capitoli.
%\part{Parte}

\chapter{Capitolo}
\section{Sezione}
\subsection{Sottosezione}
\subsubsection{Sottosottosezione}

\chapter{Background}

\chapter{Validation metrics}
\label{ch:metrics}

% =============================================================================
% Convention used in this file:
%   - every metric has: a 2-3 line description, "Inputs", "Range", "Limits".
%   - "% CITE:" comments list the papers to cite (BibTeX keys in thud.bib).
%     Keys marked [Paper/] are in the project's Paper folder; keys marked
%     [added] were appended to thud.bib and are NOT in Paper/ (verify them).
%   - "% TODO:" comments mark what to fill in with project data.
% =============================================================================

\section{Introduction and motivation}
\label{sec:metrics-intro}
% Why evaluating generated documentation is hard: no single ground truth,
% many valid paraphrases, but small wording changes ("must" vs "must not")
% can invert the contract.
% CITE: rai_review_2022 [Paper/]  (quality attributes of code documentation)
% CITE: zhang_survey_2022, zhu_automatic_2019 [Paper/]  (how summarization is evaluated)
% CITE: hoang_codewiki_2026 [Paper/]  (BLEU/ROUGE fail to capture documentation quality)
% TODO: state the research question answered by this chapter and by Chapter~\ref{ch:validation}.

\section{Evaluation setting and notation}
\label{sec:metrics-setting}
% Define: function $f$, AST metadata $\mathcal{A}(f)$, generated documentation $d_{\mathrm{gen}}$,
% reference documentation $d_{\mathrm{ref}}$ (ground truth, dataset/ground_truth.jsonl),
% source code $c$.
% Define the three kinds of input pairs used below:
%   [AST vs Doc], [Doc vs Doc], [Doc vs Code / behaviour].
% TODO: short description of the benchmark dataset (libraries, number of functions).

\section{Metric taxonomy}
\label{sec:metrics-taxonomy}
% Overview figure/table: 4 levels, input pair, needs reference?, needs LLM?, deterministic?
% CITE: yang_docagent_2025 [Paper/]  (Completeness / Helpfulness / Truthfulness axes)
% CITE: zhang_survey_2022, zhu_automatic_2019 [Paper/]  (classic evaluation families)
\begin{table}[ht]
	\centering
	\caption{Metric levels used in this thesis.}
	\label{tab:metric-taxonomy}
	\begin{tabular}{lll}
		\hline
		\textbf{Level} & \textbf{What it measures} & \textbf{Reference needed?} \\
		\hline
		I\ \ Formal / AST        & Contract consistency with the compiler AST  & No (AST is the oracle) \\
		II\ Software quality     & Actionability, error and edge-case coverage & No (source code is the oracle) \\
		III Semantic similarity  & Closeness to reference documentation       & Yes \\
		III-b LLM-as-a-judge     & Faithfulness and alignment, graded 1--5    & Optional \\
		IV\ Downstream           & Retrieval utility and behavioural fidelity & No / tests \\
		\hline
	\end{tabular}
\end{table}

% -----------------------------------------------------------------------------
\section{Level I: Formal and AST-based metrics}
\label{sec:metrics-level1}
% Common idea: libclang AST is a deterministic oracle for signature and contract.
% CITE (whole level): pomian_together_2024 [Paper/]  (static analysis as a guard on LLM output)
% CITE (whole level): lomshakov_proconsul_2024 [Paper/]  (compiler-provided structure for summarization)

\subsection{Verifier Pass Rate}
\label{sec:m-verifier}
% Fraction of functions whose documentation passes all formal checks of the Verifier
% (existence ratio $\ge 0.80$, tags consistent with the AST, mandatory tags present).
% Inputs: [AST vs Doc]. Range: [0,100]\%. Limit: checks form, not semantics.
% CITE: yang_docagent_2025 [Paper/]  (deterministic verifier + Truthfulness)
% CITE: pomian_together_2024 [Paper/]

\subsection{Parameter Precision, Recall and F1}
\label{sec:m-param-f1}
% Compares the set of \texttt{@param} names with the formal parameters of the AST
% (slot filling). Precision penalises invented parameters, recall penalises omitted ones.
% Inputs: [AST vs Doc]. Range: [0,1]. Limit: says nothing about the quality of each description.
% CITE: yang_docagent_2025 [Paper/]  (Completeness: documented parameters vs. code)
% CITE: rai_review_2022 [Paper/]  (completeness as a documentation-quality attribute)

\subsection{Return Contract Match}
\label{sec:m-return}
% Binary check: \texttt{void} functions must not document \texttt{@return};
% non-void functions must.
% Inputs: [AST vs Doc]. Range: \{0,1\}. Limit: presence only, not correctness of the description.
% CITE: yang_docagent_2025 [Paper/]

\subsection{Existence Ratio and Hallucination Rate}
\label{sec:m-hallucination}
% Existence Ratio: share of code entities named in the documentation that exist in the
% call graph / standard library. Hallucination Rate: complementary count of invented symbols.
% Inputs: [AST/call graph vs Doc]. Range: [0,1] and [0,100]\%.
% Limit: detects invented names, not wrong claims about real entities.
% CITE: yang_docagent_2025 [Paper/]  (introduces Existence Ratio for Truthfulness)
% CITE: ji_survey_2023 [added]  (definition and taxonomy of hallucination in NLG)
% CITE: lomshakov_proconsul_2024 [Paper/]  (hallucinations reduced with project context)

% -----------------------------------------------------------------------------
\section{Level II: Software-quality metrics}
\label{sec:metrics-level2}
% Metrics designed in this thesis; they have no direct counterpart in the literature, so
% the justification is the quality attributes of documentation and the DocAgent axes.
% CITE (whole level): rai_review_2022 [Paper/]  (accuracy, completeness, ... of documentation)
% CITE (whole level): yang_docagent_2025 [Paper/]  (Helpfulness axis)

\subsection{Actionability Score}
\label{sec:m-actionability}
% Weighted sum of four binary/continuous signals:
% \[ \mathrm{AS} = 0.35\,S_{\mathrm{dir}} + 0.20\,S_{\mathrm{brief}} + 0.25\,S_{\mathrm{ret}} + 0.20\,S_{\mathrm{mem}} \]
% directionality of parameters ([in]/[out]/[in,out]), quality of the brief, return clause,
% memory-safety clauses (ownership, \texttt{@pre}, \texttt{@warning}).
% Inputs: [Doc (+AST)]. Range: [0,1]. Limit: weights are heuristic -> justify with ablation.
% CITE: yang_docagent_2025 [Paper/]  (Helpfulness)
% CITE: rai_review_2022 [Paper/]
% CITE: khan_automatic_2023 [Paper/]  (structured docstring templates with @param/@return)

\subsection{Error Documentation Rate}
\label{sec:m-edr}
% Share of error branches present in the source (early returns, error codes, NULL returns)
% that are described in the documentation.
% Inputs: [Code/AST vs Doc]. Range: [0,100]\%.
% CITE: rai_review_2022 [Paper/]  (completeness)
% CITE: yang_docagent_2025 [Paper/]

\subsection{Edge Case Coverage}
\label{sec:m-ecc}
% Coverage of the guards found in the code (null, zero/negative, empty/boundary)
% by the corresponding statements in the documentation.
% Inputs: [Code/AST vs Doc]. Range: [0,100]\%.
% CITE: rai_review_2022 [Paper/]
% CITE: claessen_quickcheck_2000 [added]  (boundary values motivate property-based testing, see \ref{sec:m-roundtrip})

\subsection{Concept Checklist Score}
\label{sec:m-checklist}
% Share of critical facts (ownership, bounds, error behaviour, side effects) extracted from the
% code that are preserved by the documentation.
% Inputs: [Code facts vs Doc]. Range: [0,1].
% CITE: hoang_codewiki_2026 [Paper/]  (rubric-based evaluation derived from reference material)
% CITE: patel_assetopsbench_2026 [Paper/]  (rubric / characteristic form as soft ground truth)

% -----------------------------------------------------------------------------
\section{Level III: Semantic and NLP metrics}
\label{sec:metrics-level3}
% Reference-based metrics: compare $d_{\mathrm{gen}}$ with $d_{\mathrm{ref}}$.
% CITE (whole level): zhang_survey_2022, zhu_deep_2024 [Paper/]  (BLEU/ROUGE/METEOR as standard in code summarization)
% CITE (whole level): wu_can_2024 [Paper/]  (compares BLEU/ROUGE/METEOR/BERTScore with an LLM evaluator for code summaries)

\subsection{Lexical metrics: BLEU, ROUGE-L, TF-IDF cosine}
\label{sec:m-lexical}
% BLEU: n-gram precision with brevity penalty. ROUGE-L: longest common subsequence.
% TF-IDF cosine: bag-of-words vectors weighted by term specificity.
% Inputs: [Doc vs Doc]. Range: [0,1]. Limit: no semantics, fail on paraphrases.
% CITE: papineni_bleu_2002, lin_rouge_2004 [Paper/ bib]
% CITE: sparckjones_statistical_1972 [added]  (TF-IDF)
% CITE: zhu_deep_2024 [Paper/]  (TF-IDF/BM25 in IR-based code summarization)
% CITE: hoang_codewiki_2026 [Paper/]  (limits of BLEU/ROUGE for documentation)

\subsection{METEOR}
\label{sec:m-meteor}
% Unigram alignment with exact, stem and WordNet-synonym matches and a fragmentation penalty.
% Inputs: [Doc vs Doc]. Range: [0,1].
% CITE: banerjee_meteor_2005 [bib]
% CITE: wu_can_2024, xu_prompt_2024, geng_large_2023 [Paper/]  (METEOR used for code summaries)

\subsection{Sentence-BERT cosine similarity}
\label{sec:m-sbert}
% Cosine similarity between mean-pooled sentence embeddings (all-MiniLM-L6-v2, 384 dim.).
% \[ \mathrm{sim}(d_{\mathrm{gen}},d_{\mathrm{ref}}) = \frac{\mathbf{e}_{\mathrm{gen}}\cdot\mathbf{e}_{\mathrm{ref}}}{\|\mathbf{e}_{\mathrm{gen}}\|\,\|\mathbf{e}_{\mathrm{ref}}\|} \]
% Inputs: [Doc vs Doc]. Limit: 256-token window; blind to negations.
% CITE: reimers_sentence-bert_2019 [bib]
% CITE: zhu_reposummary_2025 [Paper/]  (embeddings to compare/cluster repository summaries)

\subsection{BERTScore and CodeBERTScore}
\label{sec:m-bertscore}
% Greedy token-to-token matching of contextual embeddings with IDF weights -> P, R, F1.
% CodeBERTScore uses a code-specific encoder (microsoft/codebert-base).
% Inputs: [Doc vs Doc] (BERTScore), [Doc vs Doc / Code] (CodeBERTScore). Limit: 512 tokens,
% score compression, insensitive to polarity.
% CITE: zhang_bertscore_2020 [bib]
% CITE: feng_codebert_2020 [Paper/]  (encoder underlying CodeBERTScore)
% CITE: guo_graphcodebert_2021 [Paper/]  (data-flow-aware variant, optional)
% CITE: zhou_codebertscore_2023 [added]  (CodeBERTScore definition; also cited in zhang_deep_2024 survey)
% CITE: wu_can_2024 [Paper/]  (BERTScore on code summarization)

\subsection{BLEURT}
\label{sec:m-bleurt}
% BERT fine-tuned on human ratings to regress a quality score for a candidate vs a reference.
% Inputs: [Doc vs Doc]. Limit: trained on general text, not on technical documentation.
% CITE: sellam_bleurt_2020 [bib]

\subsection{Fr\'echet Embedding Distance}
\label{sec:m-fid}
% Distance between the Gaussian fits of the embedding sets of generated and reference
% documentation (distribution level, not per function):
% \[ d^2 = \|\mu_1-\mu_2\|^2 + \mathrm{tr}\bigl(\Sigma_1+\Sigma_2-2(\Sigma_1\Sigma_2)^{1/2}\bigr) \]
% CITE: dowson_frechet_1982 [bib]  (closed form of the Fr\'echet distance)
% CITE: heusel_gans_2017 [added]  (FID: use on embedding distributions)

\subsection{LLM-as-a-Judge}
\label{sec:m-judge}
% An LLM grades \emph{Faithfulness} (agreement with the code) and \emph{Alignment}
% (agreement with the reference/intent) on a 1--5 scale using a rubric with a defect taxonomy
% (DEF-1\ldots DEF-5, ALIGN-1\ldots ALIGN-4). Monte Carlo sampling: 5 runs, $T=0.4$,
% report mean $\pm$ std. Judge critique is fed back to the Writer (up to 3 iterations).
% Inputs: [Doc vs Code], [Doc vs Doc]. Limits: cost, self-preference, fluff camouflage.
% CITE: liu_g-eval_2023 [Paper/]  (CoT + form-filling rubric; probability-weighted scoring)
% CITE: zhu_judgelm_2025 [Paper/]  (judges, position/format biases, rubric-based scoring)
% CITE: bai_mt-bench-101_2024 [Paper/]  (fine-grained, multi-dimensional judging)
% CITE: wu_can_2024 [Paper/]  (LLM evaluators for code summarization -- closest to our setting)
% CITE: patel_assetopsbench_2026, hoang_codewiki_2026 [Paper/]  (rubric-based LLM-judge in agent / repository documentation benchmarks)
% CITE: yang_docagent_2025 [Paper/]  (LLM-as-judge for Completeness and Helpfulness)

% -----------------------------------------------------------------------------
\section{Level IV: Downstream and behavioural metrics}
\label{sec:metrics-level4}
% Idea: evaluate documentation by what it enables, not by how it reads.

\subsection{Code retrieval: MRR and Hit@K}
\label{sec:m-retrieval}
% Use the documentation as a query over the function database; rank of the target function.
% \[ \mathrm{MRR}=\frac{1}{|Q|}\sum_{i=1}^{|Q|}\frac{1}{\mathrm{rank}_i} \]
% Hit@K: share of queries whose target is in the top $K$.
% Inputs: [Doc vs function database]. Range: [0,1].
% CITE: iyer_summarizing_2016 [Paper/]  (MRR for code retrieval with natural-language queries)
% CITE: feng_codebert_2020, guo_graphcodebert_2021 [Paper/]  (NL code search evaluated with MRR)
% CITE: husain_codesearchnet_2019 [added]  (code search benchmark)
% CITE: voorhees_trec8_1999 [added]  (original MRR definition)

\subsection{Round-trip differential testing}
\label{sec:m-roundtrip}
% (1) A coder LLM re-implements the function \emph{from the documentation only} ($f_{\mathrm{doc}}$);
% (2) the original C/C++ is transpiled to Python ($f_{\mathrm{ref}}$);
% (3) a black-box test suite (pytest + Hypothesis) is generated and run on both.
%   Self-Consistency Pass Rate: tests passed by $f_{\mathrm{doc}}$.
%   Dual Agreement Rate: tests where $f_{\mathrm{doc}}(x)\equiv f_{\mathrm{ref}}(x)$.
% Failure taxonomy: behavioural, interface mismatch, missing symbol, other.
% Limits: stateful code and information hiding (legitimate divergence), cost of test generation.
% CITE: chen_evaluating_2021 [added]  (functional correctness via tests, pass@k)
% CITE: nijkamp_codegen_2023, xu_systematic_2022 [Paper/]  (pass@k on HumanEval for code generation)
% CITE: zhang_deep_2024 [Paper/]  (survey of evaluation of code generation)
% CITE: mckeeman_differential_1998 [added]  (differential testing)
% CITE: claessen_quickcheck_2000, maciver_hypothesis_2019 [added]  (property-based testing, Hypothesis)
% CITE: yang_swe-agent_2024 [Paper/]  (test execution as feedback for agents)

% -----------------------------------------------------------------------------
\section{Summary of the metric suite}
\label{sec:metrics-summary}
% One-page table: metric | level | input pair | range | cost | expected weakness.
% The \emph{measured} reliability of each metric goes in Chapter~\ref{ch:validation}.
% TODO: build the table once the validation chapter is final.


\chapter{Dataset}
\label{ch:dataset}

\chapter{Validation delle metrihce}
\label{ch:validation}

\include{Pipeline}






%% Fine dei capitoli normali, inizio dei capitoli-appendice (opzionali)
\appendix

%\part{Appendici}

\chapter{Titolo della prima appendice}
Sed purus libero, vestibulum ut nibh vitae, mollis ultricies augue. Pellentesque velit libero, tempor sed pulvinar non, fermentum eu leo. Duis posuere eleifend nulla eget sagittis. Nam laoreet accumsan rutrum. Interdum et malesuada fames ac ante ipsum primis in faucibus. Curabitur eget libero quis leo porttitor vehicula eget nec odio. Proin euismod interdum ligula non ultricies. Maecenas sit amet accumsan sapien.

%% Parte conclusiva del documento; tipicamente per riassunto, bibliografia e/o indice analitico.
\backmatter

%% Riassunto (opzionale)
%\summary
%Maecenas tempor elit sed arcu commodo, dapibus sagittis leo egestas. Praesent at ultrices urna. Integer et nibh in augue mollis facilisis sit amet eget magna. Fusce at porttitor sapien. Phasellus imperdiet, felis et molestie vulputate, mauris sapien tincidunt justo, in lacinia velit nisi nec ipsum. Duis elementum pharetra lorem, ut pellentesque nulla congue et. Sed eu venenatis tellus, pharetra cursus felis. Sed et luctus nunc. Aenean commodo, neque a aliquam bibendum, mauris augue fringilla justo, et scelerisque odio mi sit amet diam. Nulla at placerat nibh, nec rutrum urna. Donec ut egestas magna. Aliquam erat volutpat. Phasellus vestibulum justo sed purus mattis, vitae lacinia magna viverra. Nulla rutrum diam dui, vel semper mi mattis ac. Vestibulum ante ipsum primis in faucibus orci luctus et ultrices posuere cubilia Curae; Donec id vestibulum lectus, eget tristique est.

%% Bibliografia (praticamente obbligatoria)
\bibliographystyle{plain_\languagename}%% Carica l'omonimo file .bst, dove \languagename è la lingua attiva.
%% Nel caso in cui si usi un file .bib (consigliato)
\bibliography{thud}
%% Nel caso di bibliografia manuale, usare l'environment thebibliography.

%% Per l'indice analitico, usare il pacchetto makeidx (o analogo).

\end{document}

--- Istruzioni per l'aggiunta di nuove lingue ---
Per ogni nuova lingua utilizzata aggiungere nel preambolo il seguente spezzone:
    \addto\captionsitalian{%
        \def\abstractname{Sommario}%
        \def\acknowledgementsname{Ringraziamenti}%
        \def\authorcontactsname{Contatti dell'autore}%
        \def\candidatename{Candidato}%
        \def\chairname{Direttore}%
        \def\conclusionsname{Conclusioni}%
        \def\cosupervisorname{Co-relatore}%
        \def\cosupervisorsname{Co-relatori}%
        \def\cyclename{Ciclo}%
        \def\datename{Anno accademico}%
        \def\indexname{Indice analitico}%
        \def\institutecontactsname{Contatti dell'Istituto}%
        \def\introductionname{Introduzione}%
        \def\prefacename{Prefazione}%
        \def\reviewername{Controrelatore}%
        \def\reviewersname{Controrelatori}%
        %% Anno accademico
        \def\shortdatename{A.A.}%
        \def\summaryname{Riassunto}%
        \def\supervisorname{Relatore}%
        \def\supervisorsname{Relatori}%
        \def\thesisname{Tesi di \expandafter\ifcase\csname thud@target\endcsname Laurea\or Laurea Magistrale\or Dottorato\fi}%
        \def\tutorname{Tutor aziendale}%
        \def\tutorsname{Tutor aziendali}%
    }
sostituendo a "italian" (nella 1a riga) il nome della lingua e traducendo le varie voci.
